% ==STYLE-BEGIN==
\documentclass[UTF8,11pt]{ctexart}
\usepackage{amsmath,amssymb,amsthm}
\usepackage[a4paper,margin=2.4cm,top=2.3cm,bottom=2.3cm]{geometry}
\usepackage{fancyhdr}
\usepackage{enumitem}
\usepackage{hyperref}
\hypersetup{hidelinks}

% ---- 版式 ----
\linespread{1.15}
\setlength{\parskip}{0.35em}
\setlength{\parindent}{1.6em}
\setlist{nosep,leftmargin=2em}

% ---- 页眉 ----
\pagestyle{fancy}
\fancyhf{}
\fancyhead[L]{\small\coursename}
\fancyhead[R]{\small\thepage}
\renewcommand{\headrulewidth}{0.4pt}

% ---- 定理环境：中文用正体，标题加粗，按小节共用一条流水号 ----
\newtheoremstyle{cnthm}{0.7em}{0.7em}{\normalfont}{}{\bfseries}{}{0.5em}{}
\theoremstyle{cnthm}
\newtheorem{thm}{定理}[section]
\newtheorem{lem}[thm]{引理}
\newtheorem{prop}[thm]{命题}
\newtheorem{cor}[thm]{推论}
\newtheorem{defn}[thm]{定义}
\newtheorem{rem}[thm]{注}
\newtheorem{exa}[thm]{例}
\newtheorem{conj}[thm]{猜想}
\renewcommand{\proofname}{\bfseries 证明}
\providecommand{\coursename}{}
% ==STYLE-END==
\renewcommand{\coursename}{解析几何\,·\,0923}
\title{几何与代数 0923}
\date{}
\begin{document}
\maketitle

\section{\ 二重外积公式}

\[
\vec a\times(\vec b\times\vec c)
=(\vec a\cdot\vec c)\vec b-(\vec a\cdot\vec b)\vec c .
\]

证明（待补）。

\section{\ 混合积（平行六面体的有向体积）}

\[
\vec a\cdot\vec b\times\vec c=\vec a\times\vec b\cdot\vec c .
\]

$\vec a,\vec b,\vec c$ 是三个空间向量，它们围成一个平行六面体，
计算它的体积 $V$：
\[
V=S\cdot h,\qquad S=|\vec a\times\vec b|,\qquad
h=\left|\vec c\cdot\frac{\vec a\times\vec b}{|\vec a\times\vec b|}\right|,
\]
\[
V=|\vec c\cdot\vec a\times\vec b| .
\]
当 $\vec a,\vec b,\vec c$ 是左／右手系时，
$\vec c\cdot\vec a\times\vec b>0\,/\,<0$。

显然
$V=|\vec a\cdot\vec b\times\vec c|
=|\vec b\cdot\vec c\times\vec a|$，于是
\[
\begin{aligned}
\vec c\cdot\vec a\times\vec b
&=\vec a\times\vec b\cdot\vec c\\
&=\vec a\cdot\vec b\times\vec c .
\end{aligned}
\]
用 $(\vec a,\vec b,\vec c)$ 表示 $\vec a,\vec b,\vec c$
的混合积，即
\[
\begin{aligned}
(\vec a,\vec b,\vec c)
&=\vec a\times\vec b\cdot\vec c\\
&=\vec a\cdot\vec b\times\vec c .
\end{aligned}
\]
还能有很多种表示，只需要保证三个向量手系相同，就是一样的。

\subsubsection*{混合积的性质}

\begin{enumerate}
\item $(\vec a,\vec b,\vec c_1+\vec c_2)
=(\vec a,\vec b,\vec c_1)+(\vec a,\vec b,\vec c_2)$。

\textbf{证明}\quad 记 $LHS=\vec a\times\vec b\cdot
(\vec c_1+\vec c_2)$，点乘的分配律即可得到 $RHS$。
\item $(k\vec a,\vec b,\vec c)=k(\vec a,\vec b,\vec c)$。
\item if $\vec a,\vec b,\vec c$ 线性相关，则
$(\vec a,\vec b,\vec c)=0$；特别地，
$(\vec a,\vec b,\vec b)=0$。
\item $(\vec b,\vec a,\vec c)=-(\vec a,\vec b,\vec c)$。
\item $(\vec a+k\vec b,\vec b,\vec c)=(\vec a,\vec b,\vec c)$。
\item $(\vec e_1,\vec e_2,\vec e_3)=1$。
\end{enumerate}

这些性质反映了行列式的基本性质。

在右手系 $[O;\vec e_1,\vec e_2,\vec e_3]$ 中，
$\vec a=(x_1,y_1,z_1)$，$\vec b=(x_2,y_2,z_2)$，
$\vec c=(x_3,y_3,z_3)$，把混合积用行列式表示有
\[
\begin{vmatrix}
x_1 & y_1 & z_1\\
x_2 & y_2 & z_2\\
x_3 & y_3 & z_3
\end{vmatrix}.
\]
可知，$\vec a,\vec b,\vec c$ 共面，当且仅当混合积为零，
当且仅当坐标行列式为零。

\subsubsection*{三阶行列式的基本性质}

我们把它写成行列式语言，就可以得到三阶行列式的基本性质。

\begin{enumerate}
\item \textbf{行列式相加}
\[
\begin{vmatrix}
a_{11} & a_{12} & a_{13}+a_{13}'\\
a_{21} & a_{22} & a_{23}+a_{23}'\\
a_{31} & a_{32} & a_{33}+a_{33}'
\end{vmatrix}
=
\begin{vmatrix}
a_{11} & a_{12} & a_{13}\\
a_{21} & a_{22} & a_{23}\\
a_{31} & a_{32} & a_{33}
\end{vmatrix}
+
\begin{vmatrix}
a_{11} & a_{12} & a_{13}'\\
a_{21} & a_{22} & a_{23}'\\
a_{31} & a_{32} & a_{33}'
\end{vmatrix}.
\]
\item \textbf{行列式一列全部乘 $k$，结果为原来的 $k$ 倍}
\[
\begin{vmatrix}
ka_{11} & a_{12} & a_{13}\\
ka_{21} & a_{22} & a_{23}\\
ka_{31} & a_{32} & a_{33}
\end{vmatrix}
=k\begin{vmatrix}
a_{11} & a_{12} & a_{13}\\
a_{21} & a_{22} & a_{23}\\
a_{31} & a_{32} & a_{33}
\end{vmatrix}.
\]
\item \textbf{两列成比例，行列式为零}
\[
\begin{vmatrix}
a_{11} & a_{12} & a_{12}\\
a_{21} & a_{22} & a_{22}\\
a_{31} & a_{32} & a_{32}
\end{vmatrix}=0 .
\]
\item \textbf{两列调换，符号改变}
\[
\begin{vmatrix}
a_{11} & a_{12} & a_{13}\\
a_{21} & a_{22} & a_{23}\\
a_{31} & a_{32} & a_{33}
\end{vmatrix}
=-
\begin{vmatrix}
a_{12} & a_{11} & a_{13}\\
a_{22} & a_{21} & a_{23}\\
a_{32} & a_{31} & a_{33}
\end{vmatrix}.
\]
\item \textbf{任意一列的 $k$ 倍加到另一列上去，行列式结果不变}
\[
\begin{vmatrix}
a_{11}+ka_{12} & a_{12} & a_{13}\\
a_{21}+ka_{22} & a_{22} & a_{23}\\
a_{31}+ka_{32} & a_{32} & a_{33}
\end{vmatrix}
=
\begin{vmatrix}
a_{11} & a_{12} & a_{13}\\
a_{21} & a_{22} & a_{23}\\
a_{31} & a_{32} & a_{33}
\end{vmatrix}.
\]
\item \textbf{上三角行列式的值是主对角的积}
\[
\begin{vmatrix}
1 & 0 & 0\\
0 & 1 & 0\\
0 & 0 & 1
\end{vmatrix}=1 .
\]
若主对角线上／下方元素全是 $0$，则称该行列式是上／下三角行列式。
\item \textbf{行列式的值与其转置相同}。

可以看出，行列式的行和列是等价的，所有列所具有的性质行也具备。
\end{enumerate}

\section{\ 例：异面直线距离}

$l_1,l_2$ 任取两点 $P_1,P_2$，把向量
$\overrightarrow{P_1P_2}$ 向公垂线方向投影，投影的模即为距离，即
\[
d=\left|
\overrightarrow{P_1P_2}\cdot
\frac{\vec v_1\times\vec v_2}{|\vec v_1\times\vec v_2|}
\right| .
\]
$|\overrightarrow{T_1T_2}|$ 是异面直线的距离。下给出证明：
\[
\overrightarrow{T_1T_2}
=\overrightarrow{T_1P_1}+\overrightarrow{P_1P_2}
+\overrightarrow{P_2T_2} .
\]
设 $\vec e$ 是 $\vec v_1\times\vec v_2$ 方向上的单位向量，
\[
\overrightarrow{T_1T_2}\cdot\vec e
=\left(\overrightarrow{T_1P_1}+\overrightarrow{P_1P_2}
+\overrightarrow{P_2T_2}\right)\cdot\vec e
=\overrightarrow{P_1P_2}\cdot\vec e,
\]
两边取模即得证。

\textbf{公垂线方程：拿两个平面取交集。}\quad
我们可以把公垂线看作是两个平面相交而成，利用法向量和平面上一点
可以分别确定这两个平面。而这两个平面的法向量 $\vec n_i$ 可以通过
$\vec v_1\times(\vec v_1\times\vec v_2)$ 来得到，另一个同理。
$L_1,L_2$ 分别经过 $P_1,P_2$，取公垂线上任意一点 $P_0$ 有
$\overrightarrow{P_0P_i}\cdot\vec n_i=0$，联立即可。

\section{\ 例：$\vec d=x\vec a+y\vec b+z\vec c$ 的系数}

$\vec a,\vec b,\vec c$ 不共面，$\vec d=x\vec a+y\vec b+z\vec c$，
求 $x,y,z$。

\textbf{思路：}两边同时乘两向量叉乘消元。由
\[
\vec d\cdot\vec b\times\vec c
=x\,\vec a\cdot\vec b\times\vec c
\]
即有
\[
x=\frac{\vec d\cdot\vec b\times\vec c}
{\vec a\cdot\vec b\times\vec c},\qquad
y=\frac{\vec a\cdot\vec d\times\vec c}
{\vec a\cdot\vec b\times\vec c},\qquad
z=\frac{\vec a\cdot\vec b\times\vec d}
{\vec a\cdot\vec b\times\vec c} .
\]
这就是\textbf{克莱姆法则}。

\end{document}
